\section{Data construction and timing}
\label{sec:data_details}
The original client code takes priority over firm attribution. Client code \texttt{0} is treated as missing. Client-identified and firm-only records remain distinct, and records lacking both identifiers are excluded; none occurs in the reported samples. Firm attribution can pool multiple clients. Execution messages are retained only if exactly one of the passive and aggressive flags is set. In the three samples, paired trade messages have either one passive and one aggressive side or two unassigned sides.

A trading partition consists of trading date, venue, market, instrument, and the exchange market mechanism code (EMM). Fill aggregation requires the same partition, participant, order identifier, role, and side. Passive fills also share a resting price; aggressive fills can span prices. Consecutive fills must be no more than $100$ milliseconds apart, with no intervening cancellation, modification, or new-order lifecycle event for the order. Original message identifiers and membership are retained. Sensitivity to this gap has not been measured.

\texttt{SEQUENCETIME} is defined as nanoseconds since the UTC epoch and is used for message ordering, book reconstruction, order age, cancellation, and subsequent states. Execution time uses \texttt{TRADETIME}. Required clocks are not replaced by another field when missing. Comparisons use integer nanoseconds. Observed trade times lie on a $1$-millisecond grid, whereas sequence times have sub-millisecond values. Selected trade rows have trade time preceding sequence time. Their numerical alignment motivates, but does not independently certify, the assumed common origin of the two clocks.

For execution $c$, let $t_c^{\mathrm{start}}$ and $t_c^{\mathrm{end}}$ be its first and last trade times, $r_c^{\mathrm{first}}$ and $r_c^{\mathrm{last}}$ the corresponding source positions, and $\rho_c^{\mathrm{start}}$ and $\rho_c^{\mathrm{end}}$ their sequence times. Saved book states follow message application. The label $t_c^{-}$ denotes the last retained participant state strictly before the first execution message, not a new execution timestamp. Its saved price need not be the exact global price immediately before that message.

For the complete-profile comparison, the target is $\rho_c^{\mathrm{end}}+H$, with $H\in\{1,2,5,10\}$ seconds. The label $t_c^{+}$ denotes the last retained state at or before that target with source index no smaller than the final execution message's index. Observations must cover the target. The profile-change measures are missing if either required state is unavailable. Candidate membership and the market-depth denominator, by contrast, use the exact reconstructed book immediately before the first execution message. This distinction between the selected participant records and the adjacent global book is maintained throughout.


When several executions are compared, $c$ always denotes the full canonical identity, including instrument, trading partition, attributed participant, and execution role. If needed, the suppressed profile indices can be restored as $\ell_{c,s}(t)$, $I_c(t)$, and $r_{c,s}(H)$; each evaluation uses the participant and instrument attached to $c$. The generic calibration objects in Appendix~\ref{sec:depth_calibration} are instrument-level objects, not execution-specific estimates.
